% Abstract section

\textbf{Background:}

Mendelian randomization (MR) has been widely used to investigate causal relationships between exposures and outcomes. However, concerns have also been raised about the reporting quality of MR studies. The STROBE-MR was developed to improve the reporting quality of MR studies, but assessments of STROBE-MR compliance are still limited and rely on manual review. With the rapid development of large language models (LLMs), it is now possible to use LLMs to automate the assessment of STROBE-MR compliance.

\textbf{Methods:}

We used OpenAI's GPT-5-mini model to assess 13,568 MR articles published between 2007 and 2025 for their compliance with the STROBE-MR checklist. Two reviewers independently evaluated the accuracy of the LLM's assessments on a random sample of 50 articles. 

\textbf{Results:}

The LLM achieved an overall accuracy of 89.0\% in assessing STROBE-MR compliance. On average, 22 of the 43 subitems of STROBE-MR were reported, with substantial variation of reporting rates across different STROBE-MR items and sections. A decreasing trend of number of reported subitems was also observed especially in recent years. 

\textbf{Conclusions:}

This study demonstrates the potential use of LLMs in automating the assessment of reporting quality. Our findings suggest that the reporting quality of MR studies needs to be improved, and more efforts are required to promote the use of STROBE-MR.

% Example cross-reference to supplementary materials
% See Table~\ref{tab:supp-example} in supplementary materials for additional details.